\section{INTRODUCTION}

%============================================================================
\subsection{Background}
\hspace{1.27cm}
The manufacturing and automation sectors have been evolving rapidly in recent years, with
robotic systems taking on a central role in improving speed, accuracy, and consistency
across many industries. One robotic design that has proven particularly effective in this
context is the Delta robot, a parallel kinematic manipulator first developed by Professor
Reymond Clavel at the \'{E}cole Polytechnique F\'{e}d\'{e}rale de Lausanne (EPFL) in 1985.
Since its introduction, the Delta robot has been widely deployed in demanding industrial
environments including food processing, pharmaceutical production, and electronics
manufacturing, where speed and precision are critical \autocite{clavel-1990}.
\par
\hspace{1.27cm}
What sets the Delta robot apart from conventional serial link manipulators is its parallel
kinematic structure, where all actuators remain mounted on a fixed base rather than along
the moving links. This keeps the end effector platform extremely lightweight, reduces
inertia, and allows the robot to achieve cycle times and repeatability levels that serial
arms simply cannot match. The mechanical design typically involves three identical kinematic
chains, each made up of a motor driven upper arm and a passive lower parallelogram linkage,
working together to position the moving platform precisely within the robot's workspace
\autocite{merlet-2006}.
\par
\hspace{1.27cm}
In parallel with these mechanical developments, the field of robot perception has advanced
considerably. The availability of affordable RGB-D cameras such as the Intel RealSense
series, combined with powerful deep learning based object detectors like YOLO, has made
real time 3D object localisation practical even outside of well funded research labs. When
paired with open source middleware like Robot Operating System 2 (ROS\,2), these tools make
it possible to build complete perception to motion pipelines that would previously have
required expensive proprietary systems.
\par
\hspace{1.27cm}
Despite these technological advances, the practical adoption of Delta robots and automated
systems more broadly remains limited in developing countries like Cambodia. High import
costs, inconsistent access to specialised hardware, and a limited local pool of engineers
trained in advanced robotics all contribute to this gap. Small manufacturers and
educational institutions in particular often lack the resources and technical support needed
to bring robotic automation into their operations.
\par
\hspace{1.27cm}
For Cambodia specifically, this gap represents a missed opportunity at a critical moment
in the country's industrial development. The food and agricultural processing sector,
which handles large volumes of rice, cassava, and packaged goods for export, still depends
heavily on manual labour for sorting and packaging tasks that Delta robots are ideally
suited to handle. Meanwhile, the Special Economic Zones around Phnom Penh and Sihanoukville
are attracting increasing foreign investment in light electronics and component assembly,
where pick and place precision and throughput directly affect competitiveness. An affordable
vision guided Delta robot system built on open source tools could offer local factories a
genuine path toward automation without the burden of expensive vendor contracts or imported
technical expertise. Beyond the factory floor, such a system also creates a concrete
platform for robotics education, helping Cambodian universities and technical institutes
train the next generation of engineers in skills such as computer vision, real time control,
and kinematics that are currently difficult to develop locally.
\par
\hspace{1.27cm}
This thesis responds to these needs by presenting the design, development, and experimental
validation of a vision based Delta robot pick and place system built on ROS\,2 using
affordable and accessible components. The system brings together a RealSense D455 depth
camera, YOLO based object detection, CAN bus motor control for three RobStride 00
actuators, and a visual end effector feedforward correction loop, all integrated within a
modular ROS\,2 architecture. The goal is to demonstrate that a fully functional vision to
motion pipeline is achievable at low cost, and to offer a replicable foundation for robotic
automation and engineering education in Cambodia and similar emerging economies.
%============================================================================
\subsection{Statement of Problem}
\hspace{1.27cm}
Commercial Delta robots offer excellent speed and accuracy but are prohibitively expensive for educational institutions and small-scale manufacturers in developing countries. Import costs, limited local technical expertise, and the absence of affordable open-source implementations create significant barriers to adoption. At the same time, the integration of computer vision into robot control introduces additional engineering challenges: coordinate frame calibration between the camera and robot, accurate 2D localization from YOLO object detection, compensation for non-linear kinematic errors across the workspace, and real-time synchronization between detection and motion planning.

\par
\hspace{1.27cm}
Existing low-cost Delta robot implementations typically rely on 2D monocular cameras with fixed-height assumptions, limiting their ability to handle objects at arbitrary heights or on moving conveyors. Furthermore, most implementations do not include closed-loop visual feedback to compensate for mechanical and calibration errors, resulting in systematic pick misses that degrade pick success rates.

\par
\hspace{1.27cm}
This thesis addresses these problems by designing a complete vision to motion pipeline that uses color detecting for 2D object detection, a calibrated rigid-body transform for camera to robot coordinate, conveyor belt for moving object picking, and a visual end-effector feedforward correction loop that iteratively minimizes the pixel-space error between the detected object and the EE marker before descent.

%============================================================================
\subsection{Goal and Objective}
\subsubsection{Goal}
\hspace{1.27cm}
The goal of this thesis is to design, develop, and validate a vision based Delta robot system capable of autonomously performing pick and place tasks on conveyor belt, using a RealSense D455 depth camera, YOLO object detection, and ROS\,2 for real-time control.

\subsubsection{Objective}
\hspace{1.27cm}
To achieve this goal, the following objectives are established:
\begin{itemize}
    \item Design and develop a functional prototype of the Delta Robot Parallel Manipulator.
    \item To visualize, generate and control trajectory for pick and place
\end{itemize}

%============================================================================
\subsection{Scope and Limitation}
\subsubsection{Scope}
\hspace{1.27cm}
The scope of this thesis includes:
\begin{itemize}
    \item Mechanical design and fabrication of the Delta robot structure using aluminum profiles and 3D-printed parts.
    \item Selection and integration of three RobStride RS-00 actuators via CAN bus at 1\,Mbit/s for each kinematic chain.
    \item Implementation of a ROS\,2 control architecture
    \item Integration of an Intel RealSense D455 depth camera for RGB object detection
    \item Implementation of a YOLO-based object detection node and full camera-to-robot coordinate transform pipeline.
    \item Visual end effector tracking using a laser dot marker and correction loop.
    \item Integration of a pneumatic solenoid gripper.
    \item Testing and demonstration of pick and place operations within the defined workspace.
\end{itemize}

\subsubsection{Limitation}
\hspace{1.27cm}
The limitations of this thesis are:
\begin{itemize}
    \item Single object class: The YOLO detector is trained for a single object class (orange blob); multi-class sorting is not implemented.
    \item No gripper feedback: The pneumatic gripper has no force or contact sensing; pick success is unverified (open-loop grasp).
    \item No GUI 
\end{itemize}
